\section{Multivariate power series in mathlib}
\label{sec:mvpowerseries}

In this section, we discuss the formalization of multivariate power series in mathlib.

A multivariate power series is defined to be a function from
\lstinline|(σ →₀ ℕ)| to \lstinline|R|, where \lstinline|(σ →₀ ℕ)| is a
finite supported function from an index type \lstinline|σ| to the natural numbers
\lstinline|ℕ| and \lstinline|R| is usually a ring.

\begin{lstlisting}
def MvPowerSeries (σ : Type*) (R : Type*) := (σ →₀ ℕ) → R
\end{lstlisting}

Mathlib provides several basic constructors for multivariate power
series.  The constructor \lstinline|MvPowerSeries.monomial|
\href{https://leanprover-community.github.io/mathlib4_docs/Mathlib/RingTheory/MvPowerSeries/Basic.html#MvPowerSeries.monomial}{\faExternalLink}
builds a single monomial term.  More precisely, if \lstinline|n : σ →₀ ℕ| is a
multi-index, then \lstinline|MvPowerSeries.monomial n| is an
\lstinline|R|-linear map from \lstinline|R| to
\lstinline|MvPowerSeries σ R|, which maps \lstinline|n| to \lstinline|1| and other to
\lstinline|0|.

Throughout this section, unless otherwise stated, all Lean code examples
are written inside the namespace \lstinline|MvPowerSeries|.

\begin{lstlisting}
noncomputable def monomial (n : σ →₀ ℕ) : R →ₗ[R] MvPowerSeries σ R :=
    LinearMap.single R (fun _ ↦ R) n
\end{lstlisting}

The notation \lstinline|→ₗ[R]| means that it is a \lstinline|R|-linear map.

The keyword \lstinline|noncomputable| means that Lean should not try to
extract executable computational content from the declaration.  This
often appears in mathlib when a construction depends on classical choice
or on definitions that are mathematically valid but not intended to be
computed by evaluation.  If many declarations in the same part of a file
are noncomputable, one can write \lstinline|noncomputable section| at the
beginning of the section, so that later declarations in that section do
not need to be individually marked with \lstinline|noncomputable def|.
For theorem proving, this is usually not a problem: one can still use
these definitions in statements and proofs in the same way as other
definitions.

A special case of monomial is the indeterminate indexed by \lstinline|s : σ|,
denoted as \lstinline|X s|.
In mathlib, \lstinline|X s|
\href{https://leanprover-community.github.io/mathlib4_docs/Mathlib/RingTheory/MvPowerSeries/Basic.html#MvPowerSeries.X}{\faExternalLink}
is defined as the monomial whose exponent
is \lstinline|1| at \lstinline|s| and \lstinline|0| elsewhere.

\begin{lstlisting}
def X (s : σ) : MvPowerSeries σ R :=
    monomial (Finsupp.single s 1) 1
\end{lstlisting}

\lstinline|Finsupp.single s 1| is a function with single support at \lstinline|s : σ|
with value \lstinline|1|.

The constant multivariate power series
\lstinline|MvPowerSeries.C|
\href{https://leanprover-community.github.io/mathlib4_docs/Mathlib/RingTheory/MvPowerSeries/Basic.html#MvPowerSeries.C}{\faExternalLink}
is defined to be the ring homomorphism from a semiring \lstinline|R| to
\lstinline|MvPowerSeries σ R|, whose underlying function is
\lstinline|monomial (0 : σ →₀ ℕ)|.


Since a multivariate power series is represented as a function from
multi-indices to coefficients, the most basic way to inspect such a
series is to extract one of its coefficients.  In mathlib this is done
using \lstinline|MvPowerSeries.coeff|
\href{https://leanprover-community.github.io/mathlib4_docs/Mathlib/RingTheory/MvPowerSeries/Basic.html#MvPowerSeries.coeff}{\faExternalLink}
.
For a multi-index
\lstinline|n : σ →₀ ℕ|, the term \lstinline|MvPowerSeries.coeff n| is an
\lstinline|R|-linear map from \lstinline|MvPowerSeries σ R| to
\lstinline|R|.  Applying it to a power series returns the coefficient of
the monomial with exponent \lstinline|n|.

\begin{lstlisting}
def coeff (n : σ →₀ ℕ) : MvPowerSeries σ R →ₗ[R] R := LinearMap.proj n
\end{lstlisting}

Since \lstinline|MvPowerSeries σ R| is definitionally the function type
\lstinline|(σ →₀ ℕ) → R|, this coefficient operation is essentially
evaluation at the multi-index \lstinline|n|.

The coefficient API gives convenient ways to describe the constructors
above.  We list below several relevant lemmas from mathlib, which express
the behavior of \lstinline|C|, \lstinline|X|, and \lstinline|monomial|
in terms of their coefficients.  The coefficient of \lstinline|C a| is
\lstinline|a| at the zero multi-index and zero elsewhere; the coefficient
of \lstinline|X s| is \lstinline|1| at the multi-index
\lstinline|Finsupp.single s 1| and zero elsewhere; and the coefficient of
\lstinline|monomial n a| at \lstinline|n| is exactly \lstinline|a|.

\begin{lstlisting}
theorem coeff_C (n : σ →₀ ℕ) (a : R) :
    coeff n (C a) = if n = 0 then a else 0

theorem coeff_X (n : σ →₀ ℕ) (s : σ) :
    coeff n (X s) = if n = Finsupp.single s 1 then 1 else 0

theorem coeff_monomial_same (n : σ →₀ ℕ) (a : R) : coeff n (monomial n a) = a
\end{lstlisting}

Coefficient lemmas are also the standard way to prove equality
of multivariate power series.  The extensionality theorem
\lstinline|MvPowerSeries.ext|
\href{https://leanprover-community.github.io/mathlib4_docs/Mathlib/RingTheory/MvPowerSeries/Basic.html#MvPowerSeries.ext}{\faExternalLink}
says that two power series are equal if all
of their coefficients are equal.

% \begin{lstlisting}
% theorem MvPowerSeries.ext
%     (h : ∀ n : σ →₀ ℕ, (MvPowerSeries.coeff n) φ = (MvPowerSeries.coeff n) ψ) :
%     φ = ψ

% theorem MvPowerSeries.ext_iff
%     {φ ψ : MvPowerSeries σ R} :
%     φ = ψ ↔
%       ∀ n : σ →₀ ℕ, (MvPowerSeries.coeff n) φ = (MvPowerSeries.coeff n) ψ
% \end{lstlisting}

Another important operation on multivariate power series is \emph{substitution}.  Suppose that \lstinline|f : MvPowerSeries σ R| is a power
series in variables indexed by \lstinline|σ|.  To substitute into
\lstinline|f|, we give, for each old variable \lstinline|s : σ|, a new
power series \lstinline|a s : MvPowerSeries τ S|.  Then
\lstinline|subst a f| should be interpreted as the result of replacing
each variable \lstinline|X s| in \lstinline|f| by the power series
\lstinline|a s|.  The resulting series has variables indexed by
\lstinline|τ|.

In mathlib, substitution is defined as a total function.  However, for an
infinite power series, not every formal substitution is mathematically
well-behaved.  Mathlib therefore introduces the predicate
\lstinline|HasSubst a|
\href{https://leanprover-community.github.io/mathlib4_docs/Mathlib/RingTheory/MvPowerSeries/Substitution.html#MvPowerSeries.HasSubst}{\faExternalLink}
, which records the conditions under which the
family \lstinline|a : σ → MvPowerSeries τ S| can be substituted into a
multivariate power series.

\begin{lstlisting}
structure HasSubst (a : σ → MvPowerSeries τ S) : Prop where
  const_coeff (s : σ) : IsNilpotent (constantCoeff (a s))
  coeff_zero (d : τ →₀ ℕ) : {s : σ | coeff d (a s) ≠ 0}.Finite
\end{lstlisting}

The first condition says that the constant coefficient of each
\lstinline|a s| is nilpotent.  The second condition says that, for each
fixed exponent \lstinline|d : τ →₀ ℕ|, finitely many of the series
\lstinline|a s| have a nonzero coefficient at \lstinline|d|.  These
conditions ensure that the coefficients arising after substitution are
finite sums.

A common special case is when the index type \lstinline|σ| is finite and
each substituted series has zero constant coefficient.  In this situation
mathlib provides the following theorem.

\begin{lstlisting}
theorem hasSubst_of_constantCoeff_zero [Finite σ] {a : σ → MvPowerSeries τ S}
    (ha : ∀ s : σ, constantCoeff (a s) = 0) : HasSubst a
\end{lstlisting}

The substitution operation itself is written as
\lstinline|MvPowerSeries.subst|
\href{https://leanprover-community.github.io/mathlib4_docs/Mathlib/RingTheory/MvPowerSeries/Substitution.html#MvPowerSeries.subst}{\faExternalLink}
.  In
the general mathlib version, the coefficient ring may change from
\lstinline|R| to an \lstinline|R|-algebra \lstinline|S|.

\begin{lstlisting}
def subst (a : σ → MvPowerSeries τ S) (f : MvPowerSeries σ R) :
    MvPowerSeries τ S
\end{lstlisting}

When the substituted family satisfies \lstinline|HasSubst a|, the
operation \lstinline|subst a| is an
\lstinline|R|-algebra homomorphism.  This algebra homomorphism is called
\lstinline|substAlgHom|
\href{https://leanprover-community.github.io/mathlib4_docs/Mathlib/RingTheory/MvPowerSeries/Substitution.html#MvPowerSeries.substAlgHom}{\faExternalLink}
in mathlib.

\begin{lstlisting}
def substAlgHom {a : σ → MvPowerSeries τ S} (ha : HasSubst a) :
    MvPowerSeries σ R →ₐ[R] MvPowerSeries τ S
\end{lstlisting}

There is also a coefficient formula for substitution.  It expresses the
coefficient of \lstinline|subst a f| at a multi-index \lstinline|e| as a
finite sum over the coefficients of the original series \lstinline|f|.

\begin{lstlisting}
theorem coeff_subst {a : σ → MvPowerSeries τ S} (ha : HasSubst a)
    (f : MvPowerSeries σ R) (e : τ →₀ ℕ) :
    coeff e (subst a f) = ∑ᶠ d : σ →₀ ℕ,
        coeff d f • coeff e (d.prod fun s n => a s ^ n)
\end{lstlisting}

Here \lstinline|∑ᶠ| denotes a finitely supported sum, called
\lstinline|finsum|
\href{https://leanprover-community.github.io/mathlib4_docs/Mathlib/Algebra/BigOperators/Finprod.html#finsum}{\faExternalLink}
in mathlib.
If a function is finitely supported, then
\lstinline|finsum| of this function is defined to be the finite sum over
its support, otherwise it is defined to be zero.

The operation \lstinline|MvPowerSeries.expand|
\href{https://leanprover-community.github.io/mathlib4_docs/Mathlib/RingTheory/MvPowerSeries/Expand.html#MvPowerSeries.expand}{\faExternalLink}
is a special case of
substitution.  Given a nonzero natural number \lstinline|p|,
\lstinline|expand p hp| replaces every variable \lstinline|X i| by
\lstinline|X i ^ p|.  In other words, it sends a monomial with exponent
\lstinline|d| to the monomial with exponent \lstinline|p • d|, where the
scalar multiplication \lstinline|p • d| multiplies every exponent in the
multi-index by \lstinline|p|.

\begin{lstlisting}
def expand (p : ℕ) (hp : p ≠ 0) :
    MvPowerSeries σ R →ₐ[R] MvPowerSeries σ R
\end{lstlisting}

The main computational lemmas for \lstinline|expand| say that constants
are unchanged, variables are sent to their \lstinline|p|-th powers, and
monomials have their multi-index scaled by \lstinline|p|.

\begin{lstlisting}
theorem expand_C (p : ℕ) (hp : p ≠ 0) (r : R) : expand p hp (C r) = C r

@[simp]
theorem expand_X (p : ℕ) (hp : p ≠ 0) (i : σ) : expand p hp (X i) = X i ^ p

@[simp]
theorem expand_monomial (p : ℕ) (hp : p ≠ 0) (d : σ →₀ ℕ) (r : R) :
    expand p hp (monomial d r) = monomial (p • d) r
\end{lstlisting}

In mathlib, \lstinline|PowerSeries|
\href{https://leanprover-community.github.io/mathlib4_docs/Mathlib/RingTheory/PowerSeries/Basic.html#PowerSeries}{\faExternalLink}
is defined to be a special case of multivariate
power series where \lstinline|σ| is \lstinline|Unit|.

For convenience, mathlib provides the command \lstinline|name_power_vars|
\href{https://leanprover-community.github.io/mathlib4_docs/Mathlib/Tactic/Ring/NamePowerVars.html#Mathlib.Tactic.elabNamePowerSeriesVariablesOver}{\faExternalLink},
which introduces local names for the variables of a multivariate power
series ring of the form \lstinline|MvPowerSeries (Fin n) R|.  Here
\lstinline|Fin n|
\href{https://leanprover-community.github.io/mathlib4_docs/Init/Prelude.html#Fin}{\faExternalLink} is the set containing \lstinline|n| natural numbers
starting at \lstinline|0|.
This is especially useful when working with a fixed small number of variables,
since one can write familiar names such as \lstinline|X|, \lstinline|Y|,
and \lstinline|Z| instead of \lstinline|MvPowerSeries.X (0 : Fin 3)|,
\lstinline|MvPowerSeries.X (1 : Fin 3)|, and
\lstinline|MvPowerSeries.X (2 : Fin 3)|.
In Lean, these notations can be introduced as follows:
\begin{lstlisting}
name_power_vars X, Y, Z over R
\end{lstlisting}

The notation introduced by this
command is local, so it is only available in the surrounding scope.